\custompart{SYSTEM SETUP}{Operating System Deployment and RHEL Workstation Initialization}
\chapter{RHEL Installation and Setup}


\section{Introduction to EDA}

Electronic Design Automation (EDA) refers to the collection of software tools, computational frameworks, and engineering workflows used in the design, simulation, verification, and physical implementation of integrated circuits and semiconductor systems. Modern semiconductor devices contain millions to billions of transistors operating at extremely small geometries, making manual design methodologies computationally impractical and operationally impossible. EDA tools therefore serve as the foundational infrastructure enabling engineers to transform conceptual circuit designs into manufacturable silicon devices. \\ 

Contemporary EDA environments encompass a wide range of specialized applications including schematic capture systems, Hardware Description Language (HDL) development environments, circuit simulators, synthesis tools, physical layout editors, parasitic extraction frameworks, timing analysis engines, and design-rule verification systems. These tools collectively support the complete semiconductor development lifecycle, from behavioral modeling and functional verification to physical layout generation and fabrication validation.\\

The increasing complexity of Very Large Scale Integration (VLSI) systems has made EDA indispensable within both academic and industrial semiconductor design workflows. Beyond improving productivity, EDA enables precision, scalability, repeatability, and verification accuracy that would otherwise be unattainable through conventional design methodologies. As a result, modern semiconductor engineering is fundamentally dependent upon robust EDA infrastructure and the computational environments that support it.\\

\section{Why Enterprise Linux Anchors Modern EDA?}

Modern EDA toolchains impose demanding requirements on the underlying operating system environment, including stability, deterministic behavior, scripting flexibility, network interoperability, multi-user access control, and long-term software compatibility. Consequently, semiconductor design organizations and research laboratories overwhelmingly deploy EDA infrastructure on Linux-based operating systems rather than general-purpose desktop platforms.\\

Linux provides a highly modular and administratively transparent environment capable of supporting complex engineering workflows involving distributed compute resources, centralized licensing systems, shared storage infrastructure, remote administration, and automated deployment methodologies. Native shell scripting capabilities, package management systems, process control utilities, and network tooling further enable scalable workstation administration and reproducible environment configuration across large laboratory deployments.\\

Among available Linux distributions, enterprise-grade platforms such as \textbf{Red Hat Enterprise Linux (RHEL)} and its derivatives have become the de facto standard within semiconductor design environments. Their widespread adoption is driven by long-term support lifecycles, predictable release stability, extensive vendor certification, and broad compatibility with commercial EDA applications. Most major semiconductor tool vendors officially validate and support their software primarily on enterprise Linux environments, making RHEL-based systems the preferred foundation for reliable and maintainable CAD infrastructure deployment.\\

Accordingly, the workstation architecture described throughout this manual is built upon a standardized RHEL-based environment designed to provide operational consistency, scalability, and long-term maintainability for semiconductor design laboratories and VLSI research workflows.\\

% ============================================================
%  Section: Red Hat Enterprise Linux as Industrial Infrastructure
% ============================================================

\section{Red Hat Enterprise Linux as Industrial Infrastructure}

The decision to standardize a semiconductor design environment on a specific Linux distribution is not a matter of administrative preference. It is an infrastructure decision with direct consequences for EDA tool compatibility, operational stability, vendor support, and long-term maintainability. Within professional semiconductor environments, Red Hat Enterprise Linux has become a preferred foundation because it provides the stability and compatibility characteristics required by commercial EDA workflows.\\

\noindent Commercial EDA tools are typically distributed as precompiled binary applications that depend on stable system libraries, predictable runtime behavior, and well-defined operating system interfaces. RHEL maintains Application Binary Interface and Application Programming Interface consistency within a major release, allowing validated tool installations to remain reliable across routine system updates. This consistency is especially important for tools that depend on components such as \texttt{glibc}, \texttt{libX11}, Motif libraries, OpenSSL, shell utilities, and kernel-level behavior.\\

\noindent RHEL also provides long-term lifecycle stability that aligns well with semiconductor design workflows. A laboratory or production design environment may depend on a specific operating system and EDA tool combination for several years. Frequent distribution-level changes can introduce dependency drift and invalidate previously working configurations. By contrast, RHEL's enterprise release model prioritizes predictable maintenance, security updates, and controlled package evolution over rapid upstream change.\\

\noindent Vendor certification is another important factor. Major EDA vendors commonly certify their tools against specific RHEL major versions and closely compatible derivatives. Operating within a certified platform reduces support risk, simplifies troubleshooting, and provides a defensible baseline when vendor escalation is required. While unsupported Linux distributions may sometimes run the same tools, they introduce additional uncertainty into installation, execution, and support workflows.\\

\noindent The RHEL ecosystem also supports maintainable administration practices through standard tools such as \texttt{dnf}, \texttt{systemd}, \texttt{lvm}, \texttt{sssd}, and \texttt{subscription-manager}. These tools provide a well-documented administrative surface for package management, service control, storage configuration, authentication integration, and system maintenance. In institutional environments where systems must be rebuilt, audited, extended, or handed over between administrators, this consistency is a practical advantage.\\

\noindent Throughout the remainder of this manual, operating system procedures, configuration examples, package names, service management commands, and administrative workflows are specified for a RHEL~9-compatible environment. Unless otherwise noted, the procedures also apply to Rocky Linux~9 and AlmaLinux~9 as binary-compatible RHEL derivatives. References to \emph{RHEL} in the operational context of this document should therefore be understood to include these derivatives unless a Red Hat subscription-specific feature is explicitly required.\\


% ============================================================
%  Section: Design Philosophy of the EDA Environment
% ============================================================
\section{Design Philosophy of the EDA Environment}

The semiconductor design infrastructure described in this manual is not assembled opportunistically.
It is engineered according to a defined set of operational principles that govern every decision from
initial workstation deployment through long-term maintenance. Understanding these principles is
necessary context for the procedures that follow, because many configuration choices that might
otherwise appear arbitrary are direct expressions of the underlying philosophy.\\

The environment is built around \textbf{standardization}. All workstations in the laboratory are
deployed from a common baseline: identical operating system version, identical package selection,
identical directory structures, and identical tool installation paths. Deviation from this baseline
is treated as an exception requiring justification, not a routine outcome of per-machine
administration. Standardization is the prerequisite for everything else. \\

From standardization follows \textbf{reproducibility}. A workstation rebuilt from documented
procedures should produce an environment functionally identical to the one it replaces. This
property matters both operationally---when a workstation fails and must be restored---and
institutionally, as the laboratory's administrative knowledge must remain embedded in its
documentation rather than in the accumulated state of individual machines.\\

The deployment is designed to \textbf{scale}. Procedures written for a single workstation apply
without modification to the full laboratory fleet. Configuration that must be individually
customized per machine is a design defect. Where centralized mechanisms exist---for user
authentication, license server connectivity, shared storage, or environment initialization---they
are used in preference to local equivalents.\\

\textbf{Centralized administration} follows from this. User accounts, access permissions, and
shell environments are managed at the infrastructure level rather than configured locally on each
workstation. This reduces the administrative surface area, eliminates inconsistency across
machines, and ensures that changes propagate uniformly rather than requiring repeated per-host
intervention.\\

Underlying all of these principles is the explicit goal of \textbf{eliminating configuration
drift}. In long-running laboratory environments, workstations tend to diverge from one another
over time as ad hoc fixes, local experiments, and undocumented modifications accumulate. The
procedures in this manual are written to resist this tendency---through documented baselines,
version-controlled configuration, and periodic validation practices that detect and correct
deviation before it compounds.\\

The result is an infrastructure in which any workstation can be administered, diagnosed, or
rebuilt by any administrator with access to this document, without dependence on institutional
memory or machine-specific knowledge.


\section{Overview of Laboratory Infrastructure}
\section{EDA Tool Requirements and System Constraints}

